\documentclass[11pt, a4paper]{article}

% Gemeinsame Style-Definitionen (TEXINPUTS muss templates/ enthalten — siehe scripts/build_tex.py)
% bfw-style.tex — gemeinsame Style-Definitionen für alle Handouts/Aufgabenhefte
% Voraussetzung: Kompilation mit xelatex (wegen fontspec)

\usepackage[ngerman]{babel}
% Babel-ngerman macht " zu einem aktiven Shorthand (z.B. für "a -> ä). In unseren
% Texten verwenden wir " als normales Zeichen — daher Shorthand komplett ausschalten.
\shorthandoff{"}
\usepackage{geometry}
\geometry{a4paper, top=2.4cm, bottom=2.4cm, left=2.3cm, right=2.3cm,
          headheight=15pt, headsep=0.7cm, footskip=1.1cm}

\usepackage{fontspec}
% Fallback-Kette: Source Sans Pro -> Calibri -> Segoe UI -> Latin Modern Sans
% (Latin Modern ist immer mit MiKTeX vorhanden)
\IfFontExistsTF{Source Sans Pro}{%
  \setmainfont{Source Sans Pro}[Ligatures=TeX]%
}{\IfFontExistsTF{Calibri}{%
  \setmainfont{Calibri}[Ligatures=TeX]%
}{\IfFontExistsTF{Segoe UI}{%
  \setmainfont{Segoe UI}[Ligatures=TeX]%
}{%
  \setmainfont{Latin Modern Sans}[Ligatures=TeX]%
}}}
\IfFontExistsTF{Source Code Pro}{%
  \setmonofont{Source Code Pro}[Scale=0.92]%
}{\IfFontExistsTF{Consolas}{%
  \setmonofont{Consolas}[Scale=0.92]%
}{%
  \setmonofont{Latin Modern Mono}[Scale=0.92]%
}}

% Gesperrte Versalien für Labels/Titelbalken (Kapitälchen-Ersatz, funktioniert
% mit jeder OpenType-Schrift — Latin Modern Sans hat keine echten Kapitälchen)
\newcommand{\bfwlabelfont}{\footnotesize\bfseries\addfontfeatures{LetterSpace=8.0}}

\usepackage{xcolor}
% Farbpalette: hell, freundlich, modern
\definecolor{bfwPrimary}{HTML}{0EA5A4}    % Petrol/Türkis - Primär
\definecolor{bfwPrimaryDark}{HTML}{0F766E} % dunkler Petrol für Headings
\definecolor{bfwAccent}{HTML}{F59E0B}     % Amber - Akzent
\definecolor{bfwSuccess}{HTML}{10B981}    % Grün - Erfolg / Lösung
\definecolor{bfwWarn}{HTML}{F97316}       % Orange - Achtung
\definecolor{bfwInk}{HTML}{0F172A}        % fast Schwarz
\definecolor{bfwInkSoft}{HTML}{475569}    % gedämpftes Grau
\definecolor{bfwBgSoft}{HTML}{F8FAFC}     % off-white
\definecolor{bfwBgTeal}{HTML}{ECFEFF}     % helles Türkis
\definecolor{bfwBgAmber}{HTML}{FEF3C7}    % helles Gelb
\definecolor{bfwBgRose}{HTML}{FFE4E6}     % helles Rosé für Eselsbrücken
\definecolor{bfwTabHead}{HTML}{E4F3F2}    % zarte Kopfzeilen-Hinterlegung für Tabellen

\usepackage[colorlinks=true, linkcolor=bfwPrimaryDark, urlcolor=bfwPrimaryDark]{hyperref}
\usepackage{lastpage}
\usepackage{enumitem}
\setlist[itemize]{leftmargin=*, itemsep=2pt, topsep=2pt}
\setlist[enumerate]{leftmargin=*, itemsep=2pt, topsep=2pt}

\usepackage[most]{tcolorbox}
\usepackage{booktabs}
\usepackage{colortbl}
\usepackage{tikz}
\usetikzlibrary{decorations.pathmorphing, positioning, shapes.geometric, arrows.meta}

% ---------------------------------------------------------------------------
% Einheitliches Tabellen-Design
%   - etwas mehr Zeilenluft, dezente graue Linien (wirkt auch mit \hline ruhiger)
%   - Kopfzeile einer Tabelle bei Bedarf zart hinterlegen: \rowcolor{bfwTabHead}
% ---------------------------------------------------------------------------
\renewcommand{\arraystretch}{1.25}
\arrayrulecolor{bfwInkSoft!50}
\setlength{\heavyrulewidth}{1pt}
\setlength{\lightrulewidth}{0.5pt}

% ---------------------------------------------------------------------------
% Boxen — klare farbige Titelbalken mit gesperrten Versal-Labels
% (Titel bleibt per Option überschreibbar: \begin{merksatz}[title={...}])
% ---------------------------------------------------------------------------
\tcbset{bfwbox/.style={%
  enhanced, breakable,
  boxrule=0.5pt, arc=2.5pt,
  left=10pt, right=10pt, top=8pt, bottom=8pt,
  toptitle=4pt, bottomtitle=4pt,
  titlerule=0pt,
  fonttitle=\bfwlabelfont,
  coltitle=white
}}

% Merkkästchen — wichtige Definition / Kernpunkt
\newtcolorbox{merksatz}[1][]{%
  bfwbox,
  colback=bfwBgTeal,
  colframe=bfwPrimaryDark,
  colbacktitle=bfwPrimaryDark,
  title={MERKE},
  #1
}

% Eselsbrücke — Mnemoniks
\newtcolorbox{eselsbruecke}[1][]{%
  bfwbox,
  colback=bfwBgRose,
  colframe=bfwAccent!85!black,
  colbacktitle=bfwAccent!85!black,
  title={ESELSBRÜCKE},
  #1
}

% Beispiel-Box
\newtcolorbox{beispiel}[1][]{%
  bfwbox,
  colback=bfwBgSoft,
  colframe=bfwInkSoft,
  colbacktitle=bfwInkSoft,
  title={BEISPIEL},
  #1
}

% Lösungs-Box (für Aufgabenhefte)
\newtcolorbox{loesung}[1][]{%
  bfwbox,
  colback=bfwSuccess!7,
  colframe=bfwSuccess!65!black,
  colbacktitle=bfwSuccess!65!black,
  title={LÖSUNG},
  #1
}

% Achtung-Box — typische Fehler / Stolperfallen
\newtcolorbox{achtung}[1][]{%
  bfwbox,
  colback=bfwWarn!6,
  colframe=bfwWarn!85!black,
  colbacktitle=bfwWarn!85!black,
  title={ACHTUNG},
  #1
}

% Formel-Box — für zentrale Formeln (groß, ruhig, ohne Titelbalken)
\newtcolorbox{formelbox}[1][]{%
  enhanced,
  colback=bfwBgTeal,
  colframe=bfwPrimaryDark,
  boxrule=1pt, arc=3pt,
  left=10pt, right=10pt, top=6pt, bottom=6pt,
  #1
}

% Glossar-Eintrag
\newcommand{\glossareintrag}[2]{%
  \par\noindent\textbf{\color{bfwPrimaryDark}#1} \enspace #2\par\smallskip
}

% ---------------------------------------------------------------------------
% Abschnitts-Design: farbiges Nummern-Quadrat + feine Linie unter \section,
% dezent farbige \subsection/\subsubsection
% ---------------------------------------------------------------------------
\usepackage{titlesec}
\newcommand{\bfwSecBadge}[1]{%
  \begin{tikzpicture}[baseline={([yshift=-0.62em]n.center)}]
    \node[fill=bfwPrimaryDark, text=white, font=\normalsize\bfseries,
          minimum width=1.55em, minimum height=1.55em, inner sep=1pt,
          rounded corners=1.5pt] (n) {#1};
  \end{tikzpicture}}
\titleformat{\section}
  {\Large\bfseries\color{bfwPrimaryDark}}
  {\bfwSecBadge{\thesection}}{0.6em}{}
  [\vspace{-0.2em}{\color{bfwPrimary!60}\titlerule[0.8pt]}]
\titleformat{\subsection}
  {\large\bfseries\color{bfwPrimaryDark}}
  {\textcolor{bfwPrimary}{\thesubsection}}{0.5em}{}
\titleformat{\subsubsection}
  {\normalsize\bfseries\color{bfwInk}}
  {\textcolor{bfwPrimary}{\thesubsubsection}}{0.4em}{}
\titlespacing*{\section}{0pt}{1.6em}{1em}
\titlespacing*{\subsection}{0pt}{1.2em}{0.5em}

% TikZ-Stil "skizzenhaft" — etwas weicher als Rough.js, aber stabil
\tikzset{
  roughbox/.style = {
    draw=bfwPrimaryDark, thick, fill=bfwBgTeal,
    rounded corners=4pt, inner sep=6pt
  },
  roughaccent/.style = {
    draw=bfwAccent, thick, fill=bfwBgAmber,
    rounded corners=4pt, inner sep=6pt
  },
  roughline/.style = {
    draw=bfwInkSoft, thick, -{Stealth[length=2.5mm]}
  }
}

% Bullet-Symbole (bewusst kein FontAwesome — vermeidet Font-Installations-Probleme)
\newcommand{\faLightbulb}{$\bullet$}
\newcommand{\faBookmark}{$\bullet$}

% ---------------------------------------------------------------------------
% Kopf-/Fußzeile
%   Kopf links:  Wortlogo „Wirtschaftsfachwirt"
%   Kopf rechts: Dokument-Kurztext, gesetzt via \kopfzeile{...}
%   Fuß links:   © Can Siebert · 2026 — Fuß rechts: Seite X von Y
%   Titelseite (Seite 1) bleibt ohne Kopf-/Fußzeile (\handouttitel setzt
%   \thispagestyle{empty}).
% ---------------------------------------------------------------------------
\usepackage{fancyhdr}
\newcommand{\bfwKopfzeileText}{}
\newcommand{\kopfzeile}[1]{\renewcommand{\bfwKopfzeileText}{#1}}
\pagestyle{fancy}
\fancyhf{}
\fancyhead[L]{\small\bfseries\color{bfwPrimaryDark}Wirtschaftsfachwirt}
\fancyhead[R]{\small\color{bfwInkSoft}\bfwKopfzeileText}
\renewcommand{\headrulewidth}{0.6pt}
\renewcommand{\headrule}{{\color{bfwPrimary}\hrule width\headwidth height 0.6pt}}
\fancyfoot[L]{\small\color{bfwInkSoft}\copyright\ Can Siebert\,\textperiodcentered\,2026}
\fancyfoot[R]{\small\color{bfwInkSoft}Seite \thepage\ von \pageref*{LastPage}}
% Auch \thispagestyle{plain} (z.B. durch \maketitle) soll leer bleiben:
\fancypagestyle{plain}{\fancyhf{}\renewcommand{\headrulewidth}{0pt}\renewcommand{\headrule}{}}

% ---------------------------------------------------------------------------
% \handouttitel{Obertitel}{Untertitel}{Kurs-Label}{Termin-Zeile}
%   Einheitlicher professioneller Titelblock: Dachzeile, farbige Headline,
%   Untertitel, farbige Trennlinie und dezente Meta-Zeile (Kurs/Termin/Dozent).
%   Ersetzt \title/\maketitle. Direkt nach \begin{document} aufrufen.
% ---------------------------------------------------------------------------
\newcommand{\handouttitel}[4]{%
  \thispagestyle{empty}%
  \begingroup
  \setlength{\parindent}{0pt}%
  {\bfwlabelfont\color{bfwPrimary}WIRTSCHAFTSFACHWIRT (IHK)\par}
  \vspace{0.55em}
  {\huge\bfseries\color{bfwPrimaryDark}#1\par}
  \vspace{0.5em}
  {\large\color{bfwInkSoft}#2\par}
  \vspace{0.9em}
  {\color{bfwPrimary}\rule{\linewidth}{2pt}}\par
  \vspace{0.55em}
  {\small\color{bfwInkSoft}#3\ \,\textperiodcentered\,\ #4\ \,\textperiodcentered\,\ Dozent: Can Siebert\par}
  \endgroup
  \vspace{1.4em}%
}


\kopfzeile{Handout VWL Lektion 3 \textperiodcentered{} Teil 1}

\begin{document}
\handouttitel{Lektion~3 \textperiodcentered{} Teil~1: Konjunktur, Geld- \& Finanzpolitik}%
  {Wie EZB und Staat auf die Konjunkturwellen reagieren}%
  {1.1 Volkswirtschaftliche Grundlagen}%
  {Sa, 29.08.2026 \textperiodcentered{} Session 1}

\begin{merksatz}[title={\bfseries Worum geht es in diesem Teil?}]
Die wirtschaftliche Entwicklung verläuft in Wellen (Konjunktur). Staat und Zentralbank
versuchen, diese Schwankungen zu glätten und die Ziele des Stabilitätsgesetzes zu erreichen.
In Teil~1 der Lektion lernen Sie den \textbf{Konjunkturzyklus} mit seinen Indikatoren, die
\textbf{Geldpolitik der EZB} (Ziel, Instrumente, Transmissionskette) und die
\textbf{Finanzpolitik} (antizyklische Fiskalpolitik, automatische Stabilisatoren,
Staatsverschuldung) kennen. In Teil~2 (Session~2, ab 10:55~Uhr) folgen die weiteren
Politikfelder, die wirtschaftspolitischen Konzeptionen sowie Außenwirtschaft und EU ---
alle Themen sind regelmäßig Gegenstand der IHK-Prüfung.
\end{merksatz}

\tableofcontents
\vspace{1em}
\hrule
\vspace{1em}

\section{Konjunktur: Zyklus und Indikatoren (Kurzüberblick)}

Jede volkswirtschaftliche Entwicklung ist Schwankungen unterworfen. Man unterscheidet
\textbf{Konjunkturschwankungen} (mehrjährige, wellenförmige Schwankungen der
gesamtwirtschaftlichen Aktivität um den langfristigen Wachstumstrend) von
\textbf{Saisonschwankungen} (regelmäßige, jahreszeitlich bedingte Schwankungen, z.\,B.
Bau im Winter, Tourismus im Sommer) und zufallsbedingten Schwankungen (z.\,B.
Naturkatastrophen, Krisen).

\subsection{Die vier Phasen des Konjunkturzyklus}

\begin{enumerate}
  \item \textbf{Aufschwung (Expansion):} steigende Nachfrage, Produktion und
    Investitionen, sinkende Arbeitslosigkeit, wachsender Optimismus.
  \item \textbf{Hochkonjunktur (Boom):} volle Kapazitätsauslastung, hohe Gewinne,
    Vollbeschäftigung, aber steigende Preise und Zinsen (Überhitzungsgefahr).
  \item \textbf{Abschwung (Rezession):} sinkende Nachfrage und Produktion,
    Auftragsrückgänge, zunehmende Arbeitslosigkeit, Pessimismus.
  \item \textbf{Tiefphase (Depression):} geringe Auslastung, hohe Arbeitslosigkeit,
    niedrige Preise/Zinsen --- Ausgangspunkt für einen neuen Aufschwung.
\end{enumerate}

\subsection{Konjunkturindikatoren}

\begin{itemize}
  \item \textbf{Frühindikatoren} (laufen der Entwicklung voraus): Auftragseingänge,
    Geschäftsklimaindex (ifo), Baugenehmigungen, Aktienkurse.
  \item \textbf{Präsenzindikatoren} (zeigen die aktuelle Lage): reales BIP,
    Kapazitätsauslastung, Industrieproduktion, Kurzarbeit, offene Stellen.
  \item \textbf{Spätindikatoren} (laufen hinterher): Arbeitslosenquote, Preisniveau
    (Inflationsrate), Insolvenzen, Lohnentwicklung.
\end{itemize}

\begin{eselsbruecke}
\textbf{F-P-S wie „Fahrplan der Konjunktur``:} \textbf{F}rüh kündigt an,
\textbf{P}räsenz zeigt an, \textbf{S}pät hinkt nach. Die Arbeitslosenquote ist ein
\emph{Spät}indikator --- Unternehmen entlassen erst, wenn der Abschwung schon läuft.
\end{eselsbruecke}

\section{Geldpolitik der Europäischen Zentralbank (EZB)}

Die Geldpolitik der EU liegt seit 1999 in den Händen der \textbf{Europäischen
Zentralbank (EZB)}. Sie steuert die Geldmengenentwicklung in der Eurozone und unterstützt
zudem die allgemeine Wirtschaftspolitik in der Gemeinschaft. \textbf{Vorrangiges Ziel}
ist nach dem Vertrag von Maastricht die \textbf{Preisniveaustabilität}: Der Anstieg des
\emph{harmonisierten Verbraucherpreisindex (HVPI)} soll mittelfristig \textbf{zwei
Prozent} betragen (seit der Strategieüberprüfung 2021 ein \emph{symmetrisches}
Inflationsziel; zuvor: „unter, aber nahe zwei Prozent``) --- damit sind sowohl Inflation
als auch Deflation ausgeschlossen. Die EZB
berücksichtigt bei ihrer Geldmengensteuerung sowohl die Preisentwicklung als auch die
Geldmengenentwicklung (\textbf{Zwei-Säulen-Strategie}).

\begin{merksatz}
Geschäftsbanken benötigen Zentralbankgeld, um Kredite vergeben zu können. Die Beschaffung
dieses Geldes heißt \textbf{Refinanzierung} --- wichtigste Refinanzierungsquelle ist die
Zentralbank. Genau hier setzen die geldpolitischen Instrumente an: Sie steuern
\emph{Liquidität} (Kreditmenge) und \emph{Zinsen} (Kreditkosten).
\end{merksatz}

\subsection{Instrument 1: Offenmarktgeschäfte}

Die Offenmarktgeschäfte stehen im Zentrum des Instrumentariums. Die EZB tritt als Käufer
oder Verkäufer von Wertpapieren am offenen Markt auf: \emph{Kauft} sie Wertpapiere, fließt
Zentralbankgeld in den Bankensektor (Geldschöpfung, Nachfrage wird belebt); \emph{verkauft}
sie, fließt Geld zurück (Geldvernichtung, Zinsen steigen, Nachfrage sinkt). Vier Kategorien:

\begin{itemize}
  \item \textbf{Hauptrefinanzierungsgeschäfte (Haupttender):} befristete Transaktionen im
    \emph{wöchentlichen} Rhythmus mit einer Laufzeit von \emph{einer Woche}
    (Pensionsgeschäfte bzw. Kredite gegen Verpfändung von Wertpapieren). Der hierfür
    festgelegte Zinssatz ist der \textbf{wichtigste Leitzins} im Euroraum.
  \item \textbf{Längerfristige Refinanzierungsgeschäfte (Basistender):} gleiche Technik,
    ursprünglich monatlich mit \emph{drei Monaten} Laufzeit (heute auch länger).
  \item \textbf{Feinsteuerungsoperationen:} gleichen kurzfristige Schwankungen der
    Bankenliquidität aus.
  \item \textbf{Strukturelle Operationen:} beeinflussen den Bedarf an Zentralbankgeld
    langfristig (z.\,B. Emission von Schuldverschreibungen).
\end{itemize}

Abgewickelt werden die Geschäfte in \textbf{Tenderverfahren} (standardisierte
Ausschreibungen): nach der Durchführung \emph{Standardtender} (24 Stunden, für Haupt- und
Basistender) oder \emph{Schnelltender} (wenige Stunden, für Feinsteuerung); nach der
Zuteilung \emph{Mengentender} (EZB gibt Volumen \emph{und} Zinssatz vor) oder
\emph{Zinstender} (Banken bieten Zinssätze, Zuteilung nach Höhe des Gebots).

\subsection{Instrument 2: Ständige Fazilitäten}

Die Fazilitäten stehen den Geschäftsbanken \emph{ständig} zur Verfügung (keine
Ausschreibung, Laufzeit: ein Geschäftstag):

\begin{itemize}
  \item \textbf{Spitzenrefinanzierungsfazilität} (Kreditfazilität): Banken können sich
    über Nacht unbegrenzt Zentralbankgeld leihen --- ihr Zinssatz bildet die
    \emph{Obergrenze} des Geldmarktzinses.
  \item \textbf{Einlagefazilität:} Banken legen überschüssige Liquidität bis zum nächsten
    Tag bei der EZB an --- ihr Zinssatz bildet die \emph{Untergrenze}.
\end{itemize}

Beide Sätze bilden zusammen den \textbf{Zinskorridor}, in dessen Mitte sich der
Hauptrefinanzierungssatz bewegt. Diese \textbf{drei Leitzinsen} (Einlagefazilität,
Hauptrefinanzierungssatz, Spitzenrefinanzierungsfazilität) müssen Sie für die Prüfung
nennen können.

\subsection{Instrument 3: Mindestreservepolitik}

Die EZB verpflichtet die Geschäftsbanken, bestimmte prozentuale Anteile der bei ihnen in
Sicht-, Termin- und Spareinlagen angelegten Gelder als \textbf{Pflichtguthaben} auf Konten
bei der Zentralbank zu unterhalten. Eine \emph{Erhöhung} des Mindestreservesatzes entzieht
den Banken Liquidität (weniger Kreditvergabe möglich), eine \emph{Senkung} wirkt umgekehrt.
Über Änderungen entscheidet der EZB-Rat.

\subsection{Expansive und restriktive Geldpolitik --- die Transmissionskette}

\begin{itemize}
  \item \textbf{Expansive Geldpolitik} („Politik des billigen Geldes``): Leitzinsen senken,
    Liquidität zuführen $\Rightarrow$ Kreditzinsen sinken $\Rightarrow$ Kreditnachfrage und
    Geldmenge steigen $\Rightarrow$ Investitionen und Konsum steigen $\Rightarrow$ BIP
    steigt, Arbeitslosigkeit sinkt --- Mittel gegen Rezession und Deflation.
  \item \textbf{Restriktive Geldpolitik}: Leitzinsen erhöhen, Liquidität verknappen
    (z.\,B. Mindestreserve erhöhen, Zuteilungsvolumen verringern) $\Rightarrow$ Zinsniveau
    steigt, Sparneigung nimmt zu $\Rightarrow$ kreditfinanzierte Investitionen und
    Konsumentenkredite gehen zurück $\Rightarrow$ gesamtwirtschaftliche Nachfrage sinkt
    $\Rightarrow$ Preisanstieg wird verringert, Konjunktur gedämpft --- Mittel gegen
    Inflation.
\end{itemize}

\begin{beispiel}
Die Inflationsrate liegt deutlich über zwei Prozent. Die EZB \emph{erhöht die Leitzinsen}:
Die Refinanzierung wird für die Geschäftsbanken teurer, sie geben die höheren Zinsen an
Unternehmen und Haushalte weiter. Kredite für Maschinen, Bauvorhaben und Konsum lohnen sich
weniger, die gesamtwirtschaftliche Nachfrage sinkt --- der Preisauftrieb lässt nach.
Kehrseite: Wachstum und Beschäftigung können leiden.
\end{beispiel}

\begin{eselsbruecke}
\textbf{„Zins runter --- Wirtschaft munter. Zins rauf --- Preisbremse drauf.``}
Expansiv = Gas geben (gegen Rezession), restriktiv = bremsen (gegen Inflation).
\end{eselsbruecke}

\subsection{Grenzen der Geldpolitik}

Die Instrumente wirken nur bedingt: Internationale Finanzströme und hohe Lohnabschlüsse
(über dem Produktivitätsfortschritt) können Inflation erzeugen; ob und wann Banken
Zinsänderungen weitergeben, ist unklar; bei pessimistischer Grundhaltung fördern selbst
niedrige Zinsen die Investitions- und Konsumbereitschaft nicht; und die EZB ist auf eine
einvernehmliche Fiskalpolitik der Regierungen angewiesen (erhöht der Staat in der
Inflationsbekämpfung gleichzeitig schuldenfinanziert seine Ausgaben, verpufft die
restriktive Wirkung).

\section{Finanzpolitik (Fiskalpolitik)}

Die Finanzpolitik umfasst alle Maßnahmen, die den \textbf{Staatshaushalt} (Einnahmen und
Ausgaben von Bund, Ländern, Gemeinden und Sozialversicherung) betreffen. Aufgaben:

\begin{itemize}
  \item \textbf{Beschaffung von Einnahmen} (v.\,a. Steuern, daneben Beiträge, Gebühren) zur
    Finanzierung öffentlicher Ausgaben; reichen sie nicht, steigt die Staatsverschuldung;
  \item \textbf{Umverteilung von Einkommen (Distribution)}, z.\,B. progressive
    Einkommensteuer, Transferleistungen, Sozialversicherung;
  \item \textbf{Stabilitätsaufgaben:} Steuerung der Konjunktur im Sinne des
    \textbf{Stabilitätsgesetzes (StabG, 1967)}.
\end{itemize}

\subsection{Antizyklische Fiskalpolitik nach Keynes}

Nach der Lehre von \textbf{John Maynard Keynes} sollen staatliche Einnahmen- und
Ausgabenpolitik dem Konjunkturverlauf \emph{entgegenwirken}:

\begin{center}
\begin{tabular}{p{3.2cm} p{5.2cm} p{5.2cm}}
\toprule
 & \textbf{Boom (restriktiv)} & \textbf{Rezession (expansiv)}\\
\midrule
Ausgaben & senken (Investitionen, Subventionen, Transfers zurückfahren) & erhöhen (öffentliche Aufträge, Konjunkturprogramme)\\
Einnahmen & erhöhen (Steuern) & senken (Steuererleichterungen)\\
Budgetwirkung & Überschüsse bilden (Rücklagen) & Defizite zulassen (\emph{Deficit Spending})\\
\bottomrule
\end{tabular}
\end{center}

Nimmt der Staat in der Rezession bewusst neue Schulden auf, um zusätzliche Nachfrage zu
erzeugen, spricht man von \textbf{Deficit Spending}.

\begin{beispiel}
In der Finanz- und Wirtschaftskrise wurden die \emph{Konjunkturpakete I und II} (2008/2009)
verabschiedet: Investitionen in Infrastruktur und öffentliche Gebäude,
Abschreibungserleichterungen (wirkten wie eine Steuererleichterung) und die
\emph{Umweltprämie (Abwrackprämie)}, die die Nachfrage nach Kraftfahrzeugen fördern sollte.
\end{beispiel}

\subsection{Automatische Stabilisatoren}

Neben gezielten Maßnahmen wirken \textbf{automatische (eingebaute) Stabilisatoren} ohne
politische Entscheidung antizyklisch: Die \emph{progressive Einkommensteuer} entzieht im
Boom überproportional Kaufkraft und entlastet in der Rezession; \emph{Arbeitslosengeld und
Sozialleistungen} stützen in der Rezession die Nachfrage, während im Aufschwung die
Beiträge die Kaufkraft dämpfen. Der Haushalt „atmet`` also mit der Konjunktur.

\subsection{Staatsverschuldung und ihre Grenzen}

Zu unterscheiden sind die jährliche \textbf{Neuverschuldung} (Kredite zum Ausgleich eines
Haushaltsdefizits) und die \textbf{Gesamtverschuldung} des Staates. Hohe Schulden und
Zinsverpflichtungen beschränken die Handlungsfähigkeit künftiger Politik. Juristische
Grenzen: die \textbf{Schuldenbremse} im Grundgesetz (Art.\,109 GG, seit 2009) und der
\emph{europäische Fiskalpakt} (2013), der die Mitgliedsländer verpflichtet, die
Schuldenbremse national zu regeln.

\section{Formeln \& Rechenbeispiele}

Die folgenden Berechnungen sind Prüfungsklassiker --- beide folgen demselben Muster:
\emph{(neu $-$ alt) durch alt mal 100}. Prägen Sie sich Formel und Rechenweg anhand der
einfachen Zahlenbeispiele ein.

\begin{merksatz}[title={\bfseries Formel 1: Wirtschaftswachstum (Veränderungsrate des realen BIP)}]
\[
\text{Wachstumsrate in \%} =
\frac{\text{BIP}_{\text{neu}} - \text{BIP}_{\text{alt}}}{\text{BIP}_{\text{alt}}} \times 100
\]
Immer mit dem \emph{realen} BIP rechnen (Preise eines Basisjahres) --- sonst misst man
Inflation statt echtem Wachstum.
\end{merksatz}

\begin{beispiel}
Das reale BIP steigt von 3.000~Mrd.~Euro (Vorjahr) auf 3.060~Mrd.~Euro.\\[0.3em]
\textbf{Schritt 1:} Veränderung berechnen: $3.060 - 3.000 = 60$~Mrd.~Euro.\\
\textbf{Schritt 2:} Durch den \emph{alten} Wert teilen: $60 \div 3.000 = 0{,}02$.\\
\textbf{Schritt 3:} Mal 100: $0{,}02 \times 100 = \textbf{2{,}0~\%}$ Wirtschaftswachstum.
\end{beispiel}

\begin{merksatz}[title={\bfseries Formel 2: Inflationsrate (Veränderungsrate des Preisindex)}]
\[
\text{Inflationsrate in \%} =
\frac{\text{Preisindex}_{\text{neu}} - \text{Preisindex}_{\text{alt}}}{\text{Preisindex}_{\text{alt}}} \times 100
\]
Grundlage ist der Verbraucherpreisindex (in der Eurozone: HVPI), der die Preisentwicklung
eines repräsentativen Warenkorbs misst.
\end{merksatz}

\begin{beispiel}
Der Verbraucherpreisindex steigt von 110 auf 115,5 Punkte.\\[0.3em]
\textbf{Schritt 1:} Veränderung: $115{,}5 - 110 = 5{,}5$ Punkte.\\
\textbf{Schritt 2:} Durch den \emph{alten} Indexstand teilen: $5{,}5 \div 110 = 0{,}05$.\\
\textbf{Schritt 3:} Mal 100: $0{,}05 \times 100 = \textbf{5{,}0~\%}$ Inflationsrate.\\[0.3em]
\emph{Typischer Fehler:} durch 100 statt durch den Basiswert 110 teilen --- das ergäbe
fälschlich 5,5~\%.
\end{beispiel}

\begin{merksatz}[title={\bfseries Wichtig für die Prüfungsvorbereitung (Teil 1)}]
Sie sollten sicher können: den \emph{Konjunkturzyklus} mit Indikatoren (F-P-S), das Ziel
und die \emph{drei Instrumente der EZB} samt \emph{Leitzinsen} und der Wirkungskette einer
Leitzinsänderung (Transmission), die Unterscheidung \emph{expansiv vs. restriktiv}, die
\emph{antizyklische Fiskalpolitik} (Boom vs. Rezession) mit automatischen Stabilisatoren
sowie die \emph{Formeln} für Wachstumsrate und Inflationsrate (Abschnitt~4). In Teil~2
folgen die weiteren Politikfelder, die Konzeptionen sowie Außenwirtschaft und EU ---
inklusive Wiederholung des gesamten Blocks~1.1.
\end{merksatz}

\vspace{1em}
\hrule
\vspace{0.8em}
\noindent\textsf{\small Quellen: Rahmenplan Wirtschaftsbezogene Qualifikationen (DIHK), Abschnitt 1.1.3.2;
Textband Volkswirtschaft (DIHK-Bildungs-GmbH), Kap.~3.2, S.~39--52. Weiter geht es in Session~2 (10:55--13:15~Uhr):
Teil~2 --- Politikfelder, Außenwirtschaft \& EU + Wiederholung Block~1.1.}

\end{document}
